\section{Introduction}
\label{sec:intro}

A correlation table is the first instrument of infrared interpretation. It lists, for each functional group, the wavenumber range in which a band is expected, together with a word or two about its intensity and shape, and it is where every chemist starts when a spectrum has to be read without a reference compound. The tables have a long lineage, from the compilations of Bellamy and of Colthup, Daly and Wiberley~\cite{bellamy1975infrared,colthup1990introduction} to the practitioner's summaries in current use, among which Smith's column in \emph{Spectroscopy}~\cite{smith2025bigreview4,smith2026bigreview8} is one of the most widely read. Their entries were derived from curated spectra of known compounds, refined by decades of experience, and they are for the most part right.

What the tables do not carry is a measure of how much each entry is worth. A rule of the form ``ester, C=O stretch, 1735 to 1755~\wn{}, strong'' says where the band is when the group is present; it does not say how often a band is found there when the group is absent, and that second number decides whether the observation is evidence. The tables could not give it, because they were built from spectra that contain the group, and a collection of positives has no denominator. The reader supplies the missing judgement from experience, and the experience is real, but it has not been written down, and it does not transfer to a student on the first day or to a program that applies the table one row at a time. We met the problem in that form, when encoding the tables as rules for automated interpretation: several entries, applied literally, disagreed with each other and with experimental spectra of compounds that unquestionably contained the group.

The gap has been felt before. The rule-based interpreters of the 1980s and 1990s, PAIRS and EXPIRS among them~\cite{woodruff1980pairs,woodruff1981rules,andreev1996expirs}, encoded correlation tables as rules with a confidence attached to each one; those confidences were set by the spectroscopist who wrote the rule, because collections large enough to measure them did not exist. More recent work trains neural networks on spectral collections to predict functional groups directly~\cite{fine2020spectral,punjabi2025chemotion} and reports the accuracy of the model, not the value of any rule a chemist could read. To our knowledge, no published correlation table has been scored rule by rule on experimental spectra of known structure, which is what this paper does.

Two things make it possible to measure the rules now. Open collections of experimental spectra with the depositor's structure attached have become available: the Chemotion repository has released several thousand ATR spectra from synthesis laboratories under an open licence~\cite{punjabi2025chemotion}, and the NIST Chemistry WebBook holds several thousand condensed-phase transmission spectra with structures~\cite{nist_webbook}. And the language for the question already exists. A correlation rule is a diagnostic test: it has a positive and a negative outcome, a target condition and, given a collection with known structures, a sensitivity, a specificity and a pair of likelihood ratios that say how much a positive or a negative result should shift the odds of the group~\cite{deeks2004likelihood}. Likelihood ratios are the natural currency because they are what a reader combines with what is already known about the sample, and because they can be compared across rules and across collections.

This paper reports such a measurement. A set of 101 rules for 60 functional groups was transcribed from one source, frozen before any confirmatory analysis, and scored on 1,539 ATR spectra from Chemotion and 5,124 transmission spectra from NIST, with the deposited structure as ground truth. We report, for every rule, its sensitivity, specificity and likelihood ratios in both collections, the effect of the intensity and shape qualifiers the source attaches, and a failure analysis that says what fills each window when the rule does not discriminate. The result is a quantified version of a table that has so far carried only positions and adjectives.

Three things the study does not do should be said at the outset. It does not propose new rules or adjust the published ones: every window and qualifier is scored as printed, and where it disagrees with the data the disagreement is reported. It does not model interpretation as an experienced reader performs it, weighing bands against one another; each rule is tested on its own, which is the form in which tables are taught and programmed. And it does not claim that any particular source is wrong. The source we used is the one we consider best documented; the method applies to any other unchanged, and the scripts and the frozen rule set are released so that it can be.
